% Generated by publication/build_sections.py; do not edit by hand.
\paragraph{Informal verdict.}
\textbf{A --- Complete solution}, with high confidence.

\paragraph{Lean coverage (partial).}
Lean proves automatic smoothness and the complete classification of polynomial solutions. The informal answer's shell construction of an infinite-dimensional family of nonpolynomial solutions is not formalized.

\subsection{Audited informal answer}
\subsubsection{1. Faithful restatement}

Let \(\beta,\lambda\in\mathbb R\), with \(|\lambda|>1\), and define
\[
F(\beta,\lambda)
=
%
\left\{
g\in C^1(\mathbb R,\mathbb R):
g'(t)=g(\lambda t)-\beta g(t)\quad\text{for every }t\in\mathbb R
\right\}.
\]
The source observes that, for every nonnegative integer \(p\), there is a nonzero polynomial of degree \(p\) in \(F(\lambda^p,\lambda)\), and asks, for every admissible pair \((\beta,\lambda)\), whether \(F(\beta,\lambda)\) contains nonpolynomial functions. There are no additional subparts or footnotes attached to Q756. 

There is no substantive typo or ambiguity. The only conventional point is whether ``natural integer'' includes \(0\); below,
\[
\mathbb N_0={0,1,2,\ldots},
\]
and the case \(p=0\) is the constant-polynomial case \(\beta=1\).

\medskip\hrule\medskip

\subsubsection{2. Final theorem}

\paragraph{Theorem}

For every \(\beta,\lambda\in\mathbb R\) satisfying \(|\lambda|>1\), the space \(F(\beta,\lambda)\) contains an infinite-dimensional vector subspace of nonzero \(C^\infty\) functions that are flat at the origin:
\[
g^{(n)}(0)=0\qquad(n\ge 0).
\]
Consequently,

\[
\boxed{\text{For every }\beta\in\mathbb R\text{ and every }|\lambda|>1,
\text{ there exist nonpolynomial functions in }F(\beta,\lambda).}
\]

Moreover, the polynomial solutions are completely classified:

\begin{itemize}
\item if \(\beta\neq \lambda^p\) for every \(p\in\mathbb N_0\), the zero function is the only polynomial solution;
\item if \(\beta=\lambda^p\), the polynomial solution space is one-dimensional and is generated by
\end{itemize}
\[
  P_p(t)
=
%
  \sum_{n=0}^{p}
  \frac{1}{n!}
  \left(\prod_{k=0}^{n-1}(\lambda^k-\lambda^p)\right)t^n,
\]
  where the empty product, for \(n=0\), equals \(1\).

The constructed nonpolynomial solutions require no boundedness or growth assumption at infinity.

\medskip\hrule\medskip

\subsection{3. Automatic regularity}

We first record that the \(C^1\) assumption automatically improves to \(C^\infty\).

\paragraph{Lemma 1}

Every \(g\in F(\beta,\lambda)\) is \(C^\infty\), and for every \(n\ge 0\),
\[
g^{(n+1)}(t)
=
%
\lambda^n g^{(n)}(\lambda t)-\beta g^{(n)}(t).
\tag{3.1}
\]

In particular,
\[
g^{(n+1)}(0)
=
%
(\lambda^n-\beta)g^{(n)}(0).
\tag{3.2}
\]

\subparagraph{Proof}

Suppose \(g\in C^m\), with \(m\ge1\). Then
\[
t\longmapsto g(\lambda t)-\beta g(t)
\]
belongs to \(C^m\). Since this function equals (g'), it follows that \(g'\in C^m\), hence \(g\in C^{m+1}\). Induction gives \(g\in C^\infty\).

Differentiating the defining equation \(n\) times gives
\[
g^{(n+1)}(t)
=
%
= \frac{d^n}{dt^n}g(\lambda t)-\beta g^{(n)}(t)
%
\lambda^n g^{(n)}(\lambda t)-\beta g^{(n)}(t).
\]
Putting \(t=0\) proves (3.2). \(\square\)

A useful consequence is immediate.

\paragraph{Corollary 2}

If \(g\in F(\beta,\lambda)\) and \(g(0)=0\), then \(g\) is flat at the origin:
\[
g^{(n)}(0)=0\qquad(n\ge0).
\]

\subparagraph{Proof}

Starting with \(g(0)=0\), recurrence (3.2) gives successively
\[
g'(0)=g''(0)=\cdots=0.
\]
\(\square\)

Thus it will suffice to construct nonzero global solutions whose two one-sided limits at \(0\) are both zero.

\medskip\hrule\medskip

\subsection{4. Polynomial solutions}

Let
\[
P(t)=\sum_{k=0}^m a_k t^k,\qquad a_m\neq0,
\]
be a polynomial solution of degree \(m\).

The coefficient of \(t^m\) in (P') is zero, whereas the coefficient of \(t^m\) in
\[
P(\lambda t)-\beta P(t)
\]
is
\[
(\lambda^m-\beta)a_m.
\]
Therefore
\[
\beta=\lambda^m.
\tag{4.1}
\]

More generally, comparison of the coefficient of \(t^k\) gives
\[
(k+1)a_{k+1}=(\lambda^k-\beta)a_k,
\qquad k\ge0.
\tag{4.2}
\]

If \(a_0=0\), recurrence (4.2) gives \(a_1=a_2=\cdots=0\), so every nonzero polynomial solution has \(a_0\neq0\).

Suppose now that \(\beta=\lambda^p\). Taking \(a_0=1\), recurrence (4.2) yields
\[
a_n
=
%
\frac{1}{n!}\prod_{k=0}^{n-1}(\lambda^k-\lambda^p).
\]
For \(n\le p\), none of the factors vanishes, because \(|\lambda|>1\) implies that the numbers \(\lambda^0,\lambda^1,\ldots\) are pairwise distinct. At the next step,
\[
a_{p+1}
=
%
= \frac{\lambda^p-\lambda^p}{p+1}a_p
%
0,
\]
and all subsequent coefficients vanish. This gives the polynomial \(P_p\) stated in the theorem.

Conversely, the preceding argument shows that any nonzero polynomial solution of degree \(m\) necessarily has \(\beta=\lambda^m\). Thus the polynomial classification is complete. \(\square\)

\medskip\hrule\medskip

\subsection{5. Extension from a fundamental annulus}

Set
\[
\rho:=|\lambda|>1
\]
and, for \(n\in\mathbb Z\), define the closed annuli
\[
A_n
=
%
\left\{
t\in\mathbb R:
\rho^n\le |t|\le\rho^{n+1}
\right\}.
\]

The fundamental annulus is
\[
A_0={t:1\le |t|\le\rho}.
\]

Let
\[
\mathcal S
=
%
C_c^\infty\bigl({t:1<|t|<\rho}\bigr).
\]
Thus every \(\varphi\in\mathcal S\) vanishes on neighborhoods of the four boundary points
\[
-\rho,\quad -1,\quad 1,\quad \rho.
\]

We show that any such seed \(\varphi\) can be extended uniquely to a smooth solution on \(\mathbb R\setminus{0}\).

\medskip\hrule\medskip

\subsubsection{5.1. Extension toward infinity}

Set
\[
G_\varphi(t)=\varphi(t),\qquad t\in A_0.
\]

Assume that \(G_\varphi\) has been defined on \(A_n\), for some \(n\ge0\). For \(x\in A_{n+1}\), define
\[
G_\varphi(x)
=
%
G_\varphi'!\left(\frac{x}{\lambda}\right)
+
\beta G_\varphi\!\left(\frac{x}{\lambda}\right).
\tag{5.1}
\]
Indeed, \(x/\lambda\in A_n\).

By construction,
\[
G_\varphi(\lambda t)
=
%
G_\varphi'(t)+\beta G_\varphi(t)
\]
on \(A_n\), which is exactly the desired equation.

Furthermore,
\[
\operatorname{supp}\bigl(G_\varphi|_{A_{n+1}}\bigr)
\subset
\lambda,\operatorname{supp}\bigl(G_\varphi|_{A_n}\bigr).
\tag{5.2}
\]
Differentiation does not enlarge support. Since the seed is supported strictly inside \(A_0\), induction shows that every outward piece is supported strictly inside its annulus. Hence adjacent pieces vanish on neighborhoods of their common boundary and glue \(C^\infty\)-smoothly.

This defines \(G_\varphi\) on \({|t|\ge1}\).

\medskip\hrule\medskip

\subsubsection{5.2. Extension toward the origin}

Suppose that \(G_\varphi\) has already been defined for
\[
|t|\ge \rho^{-m},
\]
where \(m\ge0\). We define it on the next inner shell
\[
A_{-(m+1)}
=
%
\left\{
t:\rho^{-(m+1)}\le |t|\le\rho^{-m}
\right\}.
\]

Fix a sign \(\sigma\in{-1,1}\), and put
\[
a=\sigma\rho^{-m}.
\]
For \(t\) in the corresponding component of \(A_{-(m+1)}\), define
\[
G_\varphi(t)
=
%
e^{-\beta(t-a)}G_\varphi(a)
+
\int_a^t e^{-\beta(t-s)}G_\varphi(\lambda s),ds.
\tag{5.3}
\]

For such \(s\),
\[
\rho^{-m}\le |\lambda s|\le\rho^{-m+1},
\]
so \(G_\varphi(\lambda s)\) is already known.

Differentiating (5.3) gives
\[
G_\varphi'(t)
=
%
-\beta G_\varphi(t)+G_\varphi(\lambda t),
\]
as required.

At \(t=a\), formula (5.3) reproduces the previously assigned value \(G_\varphi(a)\). The old and new pieces solve the same first-order inhomogeneous differential equation near the boundary, with the same boundary value. Since the forcing
\[
t\longmapsto G_\varphi(\lambda t)
\]
is already \(C^\infty\) there, the pieces glue \(C^\infty\)-smoothly.

At the initial boundaries \(t=\pm1\), the seed and the forcing both vanish in a neighborhood, so the same argument applies.

Induction defines a function
\[
G_\varphi\in C^\infty(\mathbb R\setminus{0})
\]
satisfying
\[
G_\varphi'(t)=G_\varphi(\lambda t)-\beta G_\varphi(t),
\qquad t\neq0.
\tag{5.4}
\]

All operations in the construction are linear in \(\varphi\).

\medskip\hrule\medskip

\subsection{6. Uniform boundedness near the origin}

We next prove that every annular extension has finite one-sided limits at \(0\).

For \(n\ge0\), set
\[
M_n
=
%
\sup_{t\in A_{-n}}|G_\varphi(t)|.
\]
Thus \(M_0=|\varphi|_\infty\).

The width of either connected component of \(A_{-n}\), for \(n\ge1\), is
\[
\ell_n
=
%
= \rho^{-n+1}-\rho^{-n}
%
(\rho-1)\rho^{-n}.
\]

Let \(t\in A_{-n}\) and let \(a\) denote the outer endpoint of its connected component. From (5.3),
\[
|G_\varphi(t)|
\le
e^{|\beta|\ell_n}|G_\varphi(a)|
+
e^{|\beta|\ell_n}
\int_a^t |G_\varphi(\lambda s)|,|ds|.
\]
Both \(G_\varphi(a)\) and \(G_\varphi(\lambda s)\) are bounded in absolute value by \(M_{n-1}\). Hence
\[
M_n
\le
e^{|\beta|\ell_n}(1+\ell_n)M_{n-1}
\le
e^{(|\beta|+1)\ell_n}M_{n-1}.
\tag{6.1}
\]

But
\[
\sum_{n=1}^\infty \ell_n
=
%
= \sum_{n=1}^\infty(\rho^{-n+1}-\rho^{-n})
%
1.
%
\]
Iterating (6.1) gives
\[
M_n
\le
e^{|\beta|+1}M_0
\qquad(n\ge1).
\tag{6.2}
\]

Thus \(G_\varphi\) is uniformly bounded on a punctured neighborhood of \(0\). Since \(G_\varphi(\lambda t)\) is also bounded there, equation (5.4) gives a uniform derivative bound:
\[
|G_\varphi'(t)|
\le
|G_\varphi(\lambda t)|+|\beta|,|G_\varphi(t)|
\le C_{\beta,\lambda,\varphi},
\qquad 0<|t|\le1.
\tag{6.3}
\]

Consequently \(G_\varphi\) is uniformly Lipschitz on each side of the origin. Therefore the finite limits
\[
L_+(\varphi)
=
%
\lim_{t\downarrow0}G_\varphi(t),
\qquad
L_-(\varphi)
=
%
\lim_{t\uparrow0}G_\varphi(t)
\tag{6.4}
\]
exist.

The map
\[
\Lambda:\mathcal S\longrightarrow\mathbb R^2,
\qquad
\Lambda(\varphi)
=
%
\bigl(L_-(\varphi),L_+(\varphi)\bigr),
\tag{6.5}
\]
is linear.

\medskip\hrule\medskip

\subsection{7. Killing both one-sided limits}

The seed space \(\mathcal S\) is infinite-dimensional. More concretely, inside the nonempty interval \((1,\rho)\) one may choose arbitrarily many pairwise disjoint intervals and place a nonzero smooth bump function in each one.

For example, on an interval \((a,b)\subset(1,\rho)\), one can take
\[
\psi_{a,b}(t)
=
%
\begin{cases}
\displaystyle
\exp\!\left(-\frac{1}{(t-a)(b-t)}\right),
& a<t<b,\\[1.2ex]
0,&\text{otherwise}.
\end{cases}
\]

Choose three linearly independent seeds
\[
\varphi_1,\varphi_2,\varphi_3\in\mathcal S.
\]
Their images under \(\Lambda\) are three vectors in the two-dimensional space \(\mathbb R^2\). Hence there exist real coefficients, not all zero, such that
\[
c_1\Lambda(\varphi_1)
+
c_2\Lambda(\varphi_2)
+
c_3\Lambda(\varphi_3)
=
%
0.
%
\]
Set
\[
\varphi=c_1\varphi_1+c_2\varphi_2+c_3\varphi_3.
\]
Since the seeds are linearly independent, \(\varphi\neq0\), while
\[
L_-(\varphi)=L_+(\varphi)=0.
\tag{7.1}
\]

Let \(G_\varphi\) be its annular extension, and define
\[
g(t)
=
%
\begin{cases}
G_\varphi(t),&t\neq0,\\
0,&t=0.
\end{cases}
\tag{7.2}
\]

By (7.1), \(g\) is continuous at \(0\). Furthermore, from the equation on the punctured line,
\[
G_\varphi'(t)
=
%
G_\varphi(\lambda t)-\beta G_\varphi(t)
\longrightarrow0
\qquad(t\to0),
\tag{7.3}
\]
because both one-sided limits of \(G_\varphi\) are zero.

For \(t\neq0\), the mean value theorem applied between \(0\) and \(t\) gives some \(\xi_t\) between \(0\) and \(t\) such that
\[
\frac{g(t)-g(0)}{t}
=
%
G_\varphi'(\xi_t).
\]
By (7.3), the right-hand side tends to zero. Thus
\[
g'(0)=0.
\]
Moreover (g') is continuous at \(0\), and the differential equation holds there as well:
\[
g'(0)=0=g(0)-\beta g(0).
\]

Therefore
\[
g\in F(\beta,\lambda).
\]

By Lemma 1, \(g\in C^\infty\). Since \(g(0)=0\), Corollary 2 gives
\[
g^{(n)}(0)=0\qquad(n\ge0).
\]
But \(g\) is nonzero, because it agrees with the nonzero seed \(\varphi\) on \(A_0\). A polynomial flat at a point must be the zero polynomial. Hence \(g\) is nonpolynomial.

This proves existence for every admissible \((\beta,\lambda)\).

\medskip\hrule\medskip

\subsubsection{Infinite-dimensionality}

The same argument gives more than one solution. Given any \(N\ge1\), choose \(N+2\) linearly independent seeds and restrict \(\Lambda\) to their (\(N+2\))-dimensional span. Its rank is at most \(2\), so its kernel has dimension at least \(N\). Since restriction to \(A_0\) recovers the seed, the corresponding global extensions remain linearly independent.

Thus the subspace of flat solutions in \(F(\beta,\lambda)\) is infinite-dimensional. \(\square\)

\medskip\hrule\medskip

\subsection{8. Analytic solutions}

The construction also explains the difference between the \(C^1\) and analytic categories.

\paragraph{Proposition}

A solution \(g\in F(\beta,\lambda)\) that is real analytic in a neighborhood of \(0\) is necessarily polynomial. More precisely:

\begin{itemize}
\item if \(\beta\neq\lambda^p\) for every \(p\in\mathbb N_0\), then \(g=0\);
\item if \(\beta=\lambda^p\), then \(g\) is a scalar multiple of \(P_p\).
\end{itemize}

\subparagraph{Proof}

From (3.2),
\[
g^{(n)}(0)
=
%
g(0)\prod_{k=0}^{n-1}(\lambda^k-\beta).
\tag{8.1}
\]

If \(g(0)=0\), all derivatives at \(0\) vanish. Analyticity gives \(g=0\) on a neighborhood of \(0\). If \(g\) vanishes on \((-\varepsilon,\varepsilon)\), then
\[
g(\lambda t)=g'(t)+\beta g(t)=0
\]
for \(|t|<\varepsilon\), so \(g\) vanishes on
\[
(-|\lambda|\varepsilon,|\lambda|\varepsilon).
\]
Iteration gives \(g=0\) on all of \(\mathbb R\).

Now suppose \(g(0)\neq0\). If no factor in (8.1) vanishes, then for all sufficiently large \(k\),
\[
|\lambda^k-\beta|\ge \frac12|\lambda|^k.
\]
Consequently the Taylor coefficient
\[
\frac{g^{(n)}(0)}{n!}
\]
grows in \(n\)-th root faster than a constant multiple of
\[
\frac{|\lambda|^{n/2}}{n},
\]
which tends to \(+\infty\). The Taylor series would therefore have radius of convergence zero, contradicting analyticity.

Hence some factor vanishes:
\[
\beta=\lambda^p.
\]
Then recurrence (8.1) terminates at order \(p\), and the Taylor series is a scalar multiple of \(P_p\). The same outward propagation argument shows equality on all of \(\mathbb R\). \(\square\)

Thus all the nonpolynomial solutions constructed above are necessarily nonanalytic at \(0\); in fact, they are nonzero and flat there.

\medskip\hrule\medskip

\subsection{9. Boundary and small-case checks}

\paragraph{The case \(p=0\)}

Here \(\beta=\lambda^0=1\), and
\[
P_0(t)=1.
\]
Indeed,
\[
P_0'(t)=0=P_0(\lambda t)-P_0(t).
\]
The theorem nevertheless supplies infinitely many nonpolynomial flat solutions in \(F(1,\lambda)\).

\paragraph{The case \(p=1\)}

Here \(\beta=\lambda\), and
\[
P_1(t)=1+(1-\lambda)t.
\]
Directly,
\[
P_1'(t)=1-\lambda
\]
and
\[
P_1(\lambda t)-\lambda P_1(t)=1-\lambda.
\]

\paragraph{Negative dilations}

When \(\lambda< -1\), the map \(t\mapsto\lambda t\) exchanges the two half-lines. This is why the construction imposes the two independent conditions
\[
L_-(\varphi)=L_+(\varphi)=0.
\]
No parity assumption and no sign assumption on \(\beta\) are needed.

\paragraph{Why the strict inequality \(|\lambda|>1\) matters}

The proof uses the geometrically shrinking annuli and the summability
\[
\sum_{n\ge1}
\left(\rho^{-n+1}-\rho^{-n}\right)
=1.
\]
The universal conclusion fails at the boundary.

If \(\lambda=1\), then
\[
g'=(1-\beta)g.
\]
For \(\beta=1\), all solutions are constant, so there are no nonpolynomial solutions.

If \(\lambda=-1\), differentiating the equation gives
\[
g''=(\beta^2-1)g,
\qquad
g'(0)=(1-\beta)g(0).
\]
For \(\beta=1\), all solutions are constant. For \(\beta=-1\), all solutions are of the form
\[
g(t)=c(1+2t),
\]
again polynomial. Thus the strict condition \(|\lambda|>1\) cannot simply be replaced by \(|\lambda|\ge1\).

\medskip\hrule\medskip

\subsubsection{Conclusion}

For every real \(\beta\) and every real \(\lambda\) with \(|\lambda|>1\), nonpolynomial solutions not only exist: the space \(F(\beta,\lambda)\) contains an infinite-dimensional family of nonzero \(C^\infty\) solutions flat at the origin. The polynomial solutions occur exactly at the resonant parameters \(\beta=\lambda^p\), but even at those parameters they form only a one-dimensional subspace inside a much larger solution space.
